\section{数据结构}
				\subsection{bitset}
					\inputminted{cpp}{sections/DataStructures/assets/DataStructure/bitset.cpp}
				\subsection{普通平衡树}
					\inputminted{cpp}{sections/DataStructures/assets/DataStructure/普通平衡树.cpp}
				\subsection{文艺平衡树}
					\inputminted{cpp}{sections/DataStructures/assets/DataStructure/文艺平衡树.cpp}
				\subsection{RangeBIT}
					\inputminted{cpp}{sections/DataStructures/assets/DataStructure/RangeBIT.cpp}
				\subsection{树状数组套权值线段树}
					\subsubsection{动态区间第 k 小}
					\inputminted{cpp}{sections/DataStructures/assets/DataStructure/动态区间第k小.cpp}
					\subsubsection{动态逆序对}
					\inputminted{cpp}{sections/DataStructures/assets/DataStructure/动态逆序对.cpp}
				\subsection{李超线段树}
					\inputminted{cpp}{sections/DataStructures/assets/DataStructure/李超线段树.cpp}
				\subsection{线段树分裂}
					\inputminted{cpp}{sections/DataStructures/assets/DataStructure/线段树分裂.cpp}
				\subsection{线段树合并}
					结构片段：需提供节点池 \texttt{node,ls,rs}，并按维护的信息补齐叶子与父节点的合并逻辑。
					\inputminted{cpp}{sections/DataStructures/assets/DataStructure/线段树合并.cpp}
				\subsection{线段树维护历史最值}
					\inputminted{cpp}{sections/DataStructures/assets/DataStructure/历史最值问题.cpp}
				\subsection{线段树套线段树}
					矩形取最大值、矩形查询最大值。
					\inputminted{cpp}{sections/DataStructures/assets/DataStructure/树套树支持矩阵取max 矩阵求max.cpp}
				\subsection{左偏树}
					\inputminted{cpp}{sections/DataStructures/assets/DataStructure/左偏树.cpp}
				\subsection{非递归线段树}
					\subsubsection{区间加，区间求最大值}
						\inputminted{cpp}{sections/DataStructures/assets/DataStructure/非递归线段树求最大值.cpp}
				\subsection{LCT 动态树}
					\subsubsection{路径最大值与直径（片段）}
					\inputminted{cpp}{sections/DataStructures/assets/DataStructure/LCT.cpp}
					\subsubsection{路径异或（完整程序）}
					\inputminted[breakanywhere=true]{cpp}{sections/DataStructures/assets/DataStructure/LCT-路径异或.cpp}

				\subsection{Hash Table}
					\inputminted[highlightlines={4,15,17,18}]{cpp}{sections/DataStructures/assets/DataStructure/hashmap.cpp}
				\subsection{基数排序}
					\inputminted{cpp}{sections/DataStructures/assets/Miscellany/RadixSort.cpp}
